\protect \section *{\nameref *{sec:plasma-waves}}
\begin{Solution}{2.{1}}
		$3k^2v_{Te}^2/\omega_{pe}^2=3k^2\lambda_{De}^2\ll1$, тобто $k\lambda_{De}\ll1$.
	
\end{Solution}
\begin{Solution}{2.{2}}
		$f_p\approx9$~МГц $>1$~МГц: хвиля відбивається; глибина проникнення $c/\omega_p\approx5$~м.
	
\end{Solution}
\begin{Solution}{2.{3}}
		$f_{\max}=f_p\approx9$~МГц при $0^\circ$ і~$f_p/\cos\theta\approx26$~МГц при $70^\circ$; для 100~МГц потрібен $\cos\theta\approx0{,}09$ ($\theta\approx85^\circ$), практично недосяжно.
	
\end{Solution}
\begin{Solution}{2.{4}}
		у~вакуумі ні: $\nabla\cdot\vect E=ikE_0\neq0$ при $\rho=0$, електромагнітні хвилі поперечні.
	У~плазмі така поздовжня хвиля існує~--- це ленгмюрівське коливання з~$\omega\approx\omega_p$.
	
\end{Solution}
\begin{Solution}{2.{5}}
		$n_e\approx2\times10^{9}~\text{см}^{-3}$.
	
\end{Solution}
\begin{Solution}{2.{6}}
		$n_e\approx2{,}4\times10^{5}~\text{см}^{-3}$, $v_\phi=c/0{,}9$, $v_g=0{,}9c$.
	
\end{Solution}
\begin{Solution}{2.{7}}
		$\Delta\varphi\approx56$~рад, $\approx9$ смуг.
	
\end{Solution}
\begin{Solution}{2.{8}}
		для 1~мм $n_{\mathrm{кр}}\approx1{,}1\times10^{15}~\text{см}^{-3}$,
	$\Delta\varphi\approx280$~рад ($\approx45$ смуг), $\delta\varphi\approx0{,}06$~рад; для $10{,}6$~мкм
	$\Delta\varphi\approx3$~рад, $\delta\varphi\approx6$~рад~--- вібрація більша за сигнал, потрібна компенсація
	(зате рефракція в~$10^4$ разів менша).
	
\end{Solution}
\begin{Solution}{2.{9}}
		$\tau=(2/c)\int_0^{x_c}dx/\sqrt{1-x/x_c}=4x_c/c$;
	$n_{\mathrm{кр}}\approx5\times10^{12}$--$7\times10^{13}~\text{см}^{-3}$.
	
\end{Solution}
